\documentclass[a4paper,11pt]{article}
\usepackage[utf8]{inputenc}
\usepackage[T1]{fontenc}
\usepackage[french]{babel}
\usepackage[left=1cm,right=1cm,top=.5cm,bottom=0cm]{geometry}
\usepackage{enumitem}
\usepackage{hyperref}
\usepackage{xcolor}
\usepackage{titlesec}
\usepackage{fontawesome5}
\usepackage{lmodern}
\usepackage[normalem]{ulem}

% --- Couleurs GE HealthCare (Violet/Magenta corporate) ---
\definecolor{primary}{RGB}{80, 35, 127}
\definecolor{secondary}{RGB}{80, 80, 80}

% --- Config Liens ---
\hypersetup{colorlinks=true, linkcolor=primary, filecolor=primary, urlcolor=primary}

% --- Titres ---
\titleformat{\section}{\large\bfseries\color{primary}\uppercase}{}{0em}{}[\titlerule]
\titlespacing{\section}{0pt}{8pt}{5pt}

\begin{document}

% --- EN-TÊTE ---
\begin{center}
    {\Large \textbf{Loïc Aron MBASSI EWOLO}} \\ \vspace{4pt}
    \textbf{\large Alternance R\&D -- IA pour la Santé / Traitement d'Images Médicales} \\
    \textit{\small Disponible dès septembre 2026} \\ \vspace{4pt}
    \small
    \faMapMarker\ Cergy, France (Mobile IDF) \quad
    \faPhone\ (+33) 7 59 52 33 98 \quad
    \faEnvelope\ \href{mailto:loicaron.mbassiewolo@ensea.fr}{loicaron.mbassiewolo@ensea.fr} \\ \vspace{2pt}
    \faLinkedin\ \href{https://www.linkedin.com/in/loic-mbassi/}{linkedin.com/in/loic-mbassi} \quad
    \faGithub\ \href{https://github.com/Nameless0l}{github.com/Nameless0l} \quad
    \faGlobe\ \href{https://mbassiloic.tech}{mbassiloic.tech}
\end{center}

\vspace{-6pt}

% --- PROFIL ---
\noindent\small{Élève-ingénieur double diplôme ENSEA (Signal, IA) / Polytechnique Yaoundé (\textbf{Mathématiques Appliquées}). Expérience en \textbf{classification d'images médicales} (ECG, collaboration Google DeepMind), traitement du signal et Deep Learning (PyTorch, TensorFlow). Stage de recherche au MILA (Institut Québécois d'IA). Participant au hackathon \textbf{UNBOXED -- Agentic AI Medicine} (Centrale Lyon / GE HealthCare, 2026).}

% --- FORMATION ---
\section{Formation}
\begin{itemize}[leftmargin=*, label=\tiny$\bullet$, nosep]
    \item \textbf{ENSEA -- École Nationale Supérieure de l'Électronique et de ses Applications} \hfill \textit{2025 -- 2027} \\
    \small{Double diplôme d'Ingénieur -- Systèmes Embarqués, Signal, IA. Cursus : Deep Learning, traitement du signal, analyse numérique.}
    \item \textbf{École Polytechnique de Yaoundé (1\textsuperscript{ère} école d'ingénieur CEMAC)} \hfill \textit{2021 -- 2025} \\
    \small{Diplôme d'Ingénieur de Conception (Master 2) : \textbf{Mathématiques Appliquées}, Algorithmique, Génie Logiciel.}
\end{itemize}

% --- COMPÉTENCES ---
\section{Compétences Techniques}
\begin{itemize}[leftmargin=*, nosep]
    \item \small{\textbf{Deep Learning / IA :} PyTorch, TensorFlow/Keras, scikit-learn, HuggingFace, classification, segmentation, détection.}
    \item \small{\textbf{Traitement d'images \& signal :} OpenCV, Mediapipe, filtrage, séries temporelles, analyse fréquentielle, reconstruction.}
    \item \small{\textbf{Maths appliquées :} Algèbre linéaire (SVD), optimisation, probabilités/statistiques, analyse numérique, modélisation.}
    \item \small{\textbf{Programmation :} Python (Expert), C/C++, MATLAB, Linux, Git, Docker.}
    \item \small{\textbf{Langues :} Français (natif), Anglais (professionnel/technique -- lecture d'articles scientifiques).}
\end{itemize}

% --- PROJETS IA SANTÉ / VISION ---
\section{Projets IA Santé \& Traitement d'Images}
\begin{itemize}[leftmargin=*, label={-}]

\item \textbf{Hackathon UNBOXED -- Agentic AI Medicine (GE HealthCare)} \hfill \textit{Mars 2026} \\
\textit{Centrale Lyon -- IA médicale} \hfill \textit{Python, Agentic AI, LLM}
\vspace{-2pt}
    \begin{itemize}[leftmargin=0.4cm, nosep, label=\tiny$\bullet$]
        \item \small{Conception d'une solution d'IA agentique pour améliorer le parcours de soins.}
        \item \small{Collaboration avec radiologues, experts IA médicale et mentors GE HealthCare (site de Buc).}
    \end{itemize}
\vspace{3pt}

\item \textbf{Projet UROP -- Classification d'Anomalies sur ECG} \hfill \textit{2024} \\
\textit{Collaboration mentors Google DeepMind} \hfill \textit{Python, TensorFlow/Keras, traitement du signal}
\vspace{-2pt}
    \begin{itemize}[leftmargin=0.4cm, nosep, label=\tiny$\bullet$]
        \item \small{Modèle de \textbf{détection et classification d'anomalies cardiaques} à partir d'images d'ECG papier.}
        \item \small{Extraction de features sur signaux physiologiques : filtrage, séries temporelles, analyse fréquentielle.}
        \item \small{Évaluation quantitative des performances (métriques de classification). Étude bibliographique préalable.}
    \end{itemize}
\vspace{3pt}

\item \textbf{Harmony of Languages -- Reconnaissance Gestuelle par Vision} \hfill \textit{2023 -- 2024} \\
\textit{Projet Académique -- Computer Vision} \hfill \textit{PyTorch, OpenCV, Mediapipe}
\vspace{-2pt}
    \begin{itemize}[leftmargin=0.4cm, nosep, label=\tiny$\bullet$]
        \item \small{Système ML de \textbf{classification de gestes} en temps réel (langue des signes camerounaise).}
        \item \small{Création d'un dataset custom annoté. Détection et segmentation des mains via Mediapipe.}
    \end{itemize}
\vspace{3pt}

\item \textbf{Violence Detection -- Classification Vidéo Temps Réel} \hfill \textit{2022} \\
\textit{Compétition -- 2\textsuperscript{ème} place} \hfill \textit{TensorFlow, OpenCV}
\vspace{-2pt}
    \begin{itemize}[leftmargin=0.4cm, nosep, label=\tiny$\bullet$]
        \item \small{Modèle de \textbf{classification vidéo en temps réel} pour la détection de violence. 2\textsuperscript{ème} place compétition nationale.}
    \end{itemize}
\vspace{3pt}

\end{itemize}

% --- RECHERCHE ---
\section{Expériences de Recherche}
\begin{itemize}[leftmargin=*, label={-}]

\item \textbf{Assistant de Recherche @ MILA (Institut Québécois d'IA)} \hfill \textit{Juin 2025 -- Sept 2025} \\
\textit{Remote -- Interprétabilité des Modèles} \hfill \textit{Python, PyTorch, TensorFlow}
\vspace{-2pt}
    \begin{itemize}[leftmargin=0.4cm, nosep, label=\tiny$\bullet$]
        \item \small{Étude du phénomène de Grokking sur des MLP : analyse bibliographique SOTA, reproduction d'expériences.}
        \item \small{Analyse mathématique de la convergence. Rédaction de rapports de synthèse pour l'équipe de recherche.}
    \end{itemize}
\vspace{3pt}

\item \textbf{Stage Recherche @ Labo ICS, Florence (Italie)} \hfill \textit{Avr 2026 -- Août 2026} \\
\textit{Neurosciences Computationnelles} \hfill \textit{Python, MATLAB, traitement du signal}
\vspace{-2pt}
    \begin{itemize}[leftmargin=0.4cm, nosep, label=\tiny$\bullet$]
        \item \small{Analyse de signaux et développement d'outils Python pour la recherche. Lecture d'articles scientifiques en anglais.}
        \item \small{Maintenance et évolution de bibliothèques MATLAB pour le laboratoire.}
    \end{itemize}
\vspace{3pt}

\end{itemize}

% --- DISTINCTIONS ---
\section{Leadership \& Distinctions}
\begin{itemize}[leftmargin=*, label=\tiny$\bullet$, nosep]
    \item \small{\textbf{1\textsuperscript{er} Prix} IndabaX Africa 2025 (IA Santé : nettoyage données médicales, déploiement API, dashboard Streamlit).}
    \item \small{\textbf{1\textsuperscript{er} Prix} CITS National Hackathon 2024 (optimisation urbaine, algorithme du voyageur de commerce, 75 équipes).}
    \item \small{\textbf{Lead AI Cell} Polytechnique Yaoundé : workshops ML/Deep Learning, collaboration Google DeepMind.}
    \item \small{\textbf{Certifications :} Supervised ML (DeepLearning.AI/Stanford), Python for Data Science (IBM/Coursera).}
\end{itemize}

\end{document}
